%% LyX 2.4.4 created this file.  For more info, see https://www.lyx.org/.
%% Do not edit unless you really know what you are doing.
\documentclass[oneside,english]{amsart}
\usepackage[T1]{fontenc}
\usepackage[latin9]{inputenc}
\usepackage{units}
\usepackage{enumitem}
\usepackage{amsthm}
\usepackage{geometry}
\geometry{verbose,tmargin=1in,bmargin=1in,lmargin=1in,rmargin=1in}

\makeatletter
%%%%%%%%%%%%%%%%%%%%%%%%%%%%%% Textclass specific LaTeX commands.
\numberwithin{equation}{section}
\numberwithin{figure}{section}
\newlength{\lyxlabelwidth}      % auxiliary length 

%%%%%%%%%%%%%%%%%%%%%%%%%%%%%% User specified LaTeX commands.
\usepackage{fancyhdr}
\usepackage{lastpage}
\pagestyle{fancy}
\usepackage{enumitem}
\usepackage{ifthen}
\everymath=\expandafter{\the\everymath\displaystyle}
%\DeclareMathSizes{10}{10}{10}{10}
\usepackage{multicol}

\usepackage{tikz}
\usetikzlibrary{patterns}
\usetikzlibrary{plotmarks}
\usepackage{pgfplots}
\pgfplotsset{compat=newest}

\usepackage{calc}
\usepackage[nomessages]{fp}

\usepackage{datetime}
% Added by lyx2lyx
\setlength{\parskip}{\smallskipamount}
\setlength{\parindent}{0pt}

\makeatother

\usepackage{babel}
\begin{document}
\renewcommand{\author}{Dr.\,James Ashe}

\newcommand{\classNum}{232}

\newcommand{\classSec}{A and B}

\newcommand{\examType}{Exam 4}

\renewcommand{\date}{Due: \formatdate{6}{5}{2013}}

%%First page's header%%%%%%%%%%%%%%%%%%%%%%
\fancypagestyle{firststyle} %%header settings for first pagr
{
	\fancyhead[L]{MTH \classNum{}(\classSec{}) \author{}\\
	\textbf{\examType{}} ~\date{}}
	\fancyhead[C]{}
	\fancyhead[R]{\textbf{Solutons\hspace{4cm}}}
	\setlength{\headsep}{14pt}
}
%%%%%%%%%%%%%%%%%%%%%%%%%%%%%%%%%%%%%%%%%%

%%Next page's header%%%%%%%%%%%%%%%%%%%%%%%
\setlist{nolistsep}
\lhead{\classNum{}(\classSec{})} 
\chead{\textbf{Name:}} 
\cfoot{\thepage{}~of~\pageref{LastPage}} 
\setlength{\headsep}{5pt} 
%%%%%%%%%%%%%%%%%%%%%%%%%%%%%%%%%%%%%%%%%%%

%%Various settings; see the preamble for other settings%%%%%
%\setenumerate[0]{leftmargin=12pt,itemindent=0pt,labelwidth=0pt}
\setlist{topsep=.5pt}
\renewcommand{\labelenumi}{\textbf{\arabic{enumi}.}}
\renewcommand{\labelenumii}{\textbf{\alph{enumii}.}}
\thispagestyle{firststyle}
%%%%%%%%%%%%%%%%%%%%%%%%%%%%%%%%%%%%%%%%%%
\newcommand{\xmax}{0}
\newcommand{\xmin}{-\xmax}
\newcommand{\xstep}{0}
\newcommand{\ymax}{0}
\newcommand{\ymin}{-\ymax}
\newcommand{\ystep}{0} 

\newcommand{\dspace}{\vspace{1cm}}

\small\textbf{Justify your answers and show your work. }Unless indicated
otherwise, each problem is worth 10 points.
\begin{enumerate}
\item ${\displaystyle \int}\cos^{3}2x\,dx$
\begin{eqnarray*}
 & = & {\displaystyle \int}\cos^{2}2x\cos2x\,dx\\
 & = & {\displaystyle \int}\left(1-\sin^{2}2x\right)\cos2x\,dx|_{u=\sin2x,\,du=2\cos2x\,dx}\\
 & = & {\displaystyle \int}(1-u^{2})\frac{\cos2x}{2\cos2x}\,dx\\
 & = & \frac{1}{2}\left(u-\frac{u^{3}}{3}\right)+C=\frac{1}{2}\sin2x+\frac{\sin^{3}2x}{6}+C
\end{eqnarray*}
\vfill
\item ${\displaystyle \int}\sin^{4}x\,dx$ (Hint:${\displaystyle \int\cos^{2}2x\,dx=\frac{x}{2}+\frac{\sin(4x)}{8}+C}$)
\begin{eqnarray*}
 & = & {\displaystyle \int\left(\frac{1-\cos2x}{2}\right)^{2}\,dx}\\
 & = & \frac{1}{4}\int1-2\cos2x+\cos^{2}2x\,dx\\
 & = & \frac{1}{4}\left(x-\sin2x+\left[\frac{x}{2}+\frac{\sin(4x)}{8}\right]\right)+C\\
 & = & \frac{x}{8}-\frac{\sin2x}{4}+\frac{\sin4x}{32}+C
\end{eqnarray*}
\vfill
\item ${\displaystyle \int\sin^{3}x}\cos x\,dx$
\begin{eqnarray*}
 & = & {\displaystyle \int}\sin^{2}x\sin x\cos x\,dx\\
 & = & {\displaystyle \int}\left(1-\cos^{2}x\right)\sin x\cos x\,dx|_{u=\cos x\,du\,-\sin x\,dx}\\
 & = & {\displaystyle \int}(1-u^{2})\frac{\sin x}{-\sin x}u\,dx={\displaystyle -\int}(u-u^{3})\,dx\\
 & = & -\left(\frac{u^{2}}{2}+\frac{u^{4}}{4}\right)+C=-\frac{\cos^{2}x}{2}+\frac{\cos^{4}x}{4}+C\\
 &  & \left(=\frac{\sin^{4}x}{4}\right)
\end{eqnarray*}
\vfill
\item ${\displaystyle \int\frac{dx}{\sqrt{25-x^{2}}}}$ Let $x=5\sin\theta$
so $dx=5\cos\theta$,
\begin{eqnarray*}
 & = & {\displaystyle \int\frac{5\cos\theta}{5\cos\theta}\,d\theta}\\
 & = & \int1\,d\theta\\
 & = & \theta+C=\sin^{-1}\left(\frac{x}{5}\right)+C
\end{eqnarray*}
\vfill
\item ${\displaystyle \int_{0}^{7}\sqrt{49-x^{2}}}\,dx$ Let $x=7\sin\theta$
so $dx=7\cos\theta$, $\sin^{-1}\left(\frac{7}{7}\right)=\frac{\pi}{2}$,
and $\sin^{-1}\left(\frac{0}{7}\right)=0$
\begin{eqnarray*}
 & = & {\displaystyle \int_{0}^{\nicefrac{\pi}{2}}7\cos x7\cos x\,dx}\\
 & = & 49{\displaystyle \int_{0}^{\nicefrac{\pi}{2}}\cos^{2}x\,dx}\\
 & = & 49\left(\frac{\theta}{2}+\frac{\sin(2\theta)}{4}\right)|_{0}^{\nicefrac{\pi}{2}}=\frac{49\pi}{4}
\end{eqnarray*}
\vfill
\item ${\displaystyle \int\frac{x^{2}}{\left(\sqrt{25-x^{2}}\right)^{3}}\,dx}$
(Hint: ${\displaystyle \int\frac{\sin^{2}\theta}{\cos^{2}\theta}\,d\theta={\displaystyle \int\tan^{2}\theta\,d\theta=\tan\theta+\theta+C)}}$
{[}5pts{]}\\
Let $x=5\sin\theta$ so $dx=5\cos\theta$, 
\begin{eqnarray*}
 & = & {\displaystyle \int\frac{25\sin^{2}\theta5\cos\theta}{\left(5\cos\theta\right)^{3}}\,d\theta}\\
 & = & \int\frac{25\sin^{2}\theta}{25\cos^{2}\theta}\,d\theta\\
 & = & {\displaystyle \int\tan^{2}\theta\,d\theta}=\tan\theta+\theta+C\\
 & = & \frac{\nicefrac{x}{5}}{\nicefrac{\sqrt{25-x^{2}}}{5}}+\sin^{-1}\left(\frac{x}{5}\right)+C=\frac{x}{\sqrt{25-x^{2}}}+\sin^{-1}\left(\frac{x}{5}\right)+C
\end{eqnarray*}
\vfill
\item ${\displaystyle \int\frac{x}{x^{2}+x-12}\,dx}$ {[}15 pts{]}\\
\begin{eqnarray*}
\frac{x}{x^{2}+x-12} & = & \frac{A}{x-3}+\frac{B}{x+4}\\
x & = & A(x+4)+B(x-3)\\
\mbox{let }x=3 & \Rightarrow & A=\frac{3}{7}\\
\mbox{let }x=-4 & \Rightarrow & B=\frac{4}{7}
\end{eqnarray*}
So, 
\begin{eqnarray*}
{\displaystyle \int\frac{x}{x^{2}+x-12}\,dx} & = & {\displaystyle \int\frac{3}{7(x-3)}+\frac{4}{7(x+4)}\,dx}\\
 & = & \frac{3}{7}\ln\left|x-3\right|+\frac{4}{7}\ln\left|x+4\right|+C\\
 & = & \ln^{\nicefrac{1}{7}}\left|(x-3)^{3}(x+4)^{4}\right|+C
\end{eqnarray*}
\vfill
\item ${\displaystyle \int}_{0}^{\ln2}4xe^{x}\,dx$ let $\begin{array}{cc}
u=4x & dv=e^{x}\,dx\\
du=4\,dx & v=e^{x}
\end{array}$so, 
\begin{eqnarray*}
 & = & {\displaystyle 4xe^{x}\Bigr|_{0}^{\ln2}-{\displaystyle 4\int_{0}^{\ln2}e^{x}\,dx}}\\
 & = & 4\left(xe^{x}-e^{x}\right)\Bigr|_{0}^{\ln2}\\
 & = & 4\left(\ln2\cdot2-2\right)-4(0-1)=8\ln2-4
\end{eqnarray*}
\vfill
\item ${\displaystyle \int_{1}^{\infty}6x^{-2}\,dx}$
\begin{eqnarray*}
 & = & {\displaystyle \lim_{b\rightarrow\infty}\int_{1}^{b}6x^{-2}\,dx}\\
 & = & \lim_{b\rightarrow\infty}-\frac{6}{x}\Bigr|_{1}^{b}=0-\left(-6\right)=6
\end{eqnarray*}
 \vfill
\item ${\displaystyle \int_{2}^{\infty}\frac{17}{x\ln x}\,dx}$
\begin{eqnarray*}
 & = & {\displaystyle \lim_{b\rightarrow\infty}\int_{1}^{b}\frac{17}{x\ln x}\,dx\Bigr|_{\begin{array}{c}
u=\ln x\\
du=\frac{dx}{x}
\end{array}}}\\
 & = & \lim_{b\rightarrow\infty}\int_{1}^{\ln b}\frac{17x}{x\cdot u}\,du=\lim_{b\rightarrow\infty}17\ln\left|\ln\left|b\right|\right|\Bigr|_{1}^{b}\\
 & = & 17\lim_{b\rightarrow\infty}\left(\ln\left|\ln b\right|-0\right)=\infty
\end{eqnarray*}
\vfill
\end{enumerate}

\end{document}
